\documentclass[11pt,a4paper]{article}
\usepackage{genesisl1-publication}
\GenesisTitle{Title of the Scientific Letter}
\GenesisShortTitle{Short letter title}
\GenesisSubtitle{Observation, hypothesis, critique, response, or concise result}
\GenesisType{Scientific Letter}
\GenesisAuthor{Author names}
\GenesisAffiliation{Affiliations or independent contributor}
\GenesisVersion{Draft 0.1}
\GenesisStatus{OPEN SCIENTIFIC COMMUNICATION}
\GenesisIdentifier{Optional identifier}
\GenesisLicense{CC BY 4.0}
\begin{document}
\MakeGenesisTitle
\begin{GenesisAbstract}
Summarise the central observation, question, argument, or result in a short paragraph.
\end{GenesisAbstract}
\section{Context}
Explain the scientific setting and why a concise letter is the appropriate form.
\CentralQuestion{What is the single question, observation, or claim the reader should examine?}
\section{Contribution}
Present the observation, hypothesis, critique, technical idea, response, or preliminary result.
\begin{GenesisEvidence}
Distinguish what was observed, measured, calculated, implemented, or documented from
what is inferred.
\end{GenesisEvidence}
\begin{GenesisInterpretation}
Explain what the evidence may mean and mention credible alternatives.
\end{GenesisInterpretation}
\begin{GenesisLimitations}
State the most important uncertainty and what evidence could revise the interpretation.
\end{GenesisLimitations}
\section{Proposed test or next step}
Optional. Describe the smallest useful experiment, analysis, replication, discussion, or
implementation step.
\begin{GenesisAction}
Invite focused critique, independent testing, data contribution, an alternative explanation,
or collaboration.
\end{GenesisAction}
\end{document}
